\begin{tikzpicture}[
    main/.style={
        rectangle,
        rounded corners=3pt,
        draw=blue!55,
        fill=blue!6,
        align=left,
        text width=9.15cm,
        minimum height=0.72cm,
        inner xsep=7pt,
        font=\sffamily\small
    },
    check/.style={
        rectangle,
        rounded corners=3pt,
        draw=blue!55,
        fill=blue!6,
        align=left,
        text width=9.15cm,
        minimum height=0.72cm,
        inner xsep=7pt,
        font=\sffamily\small
    },
    arrow/.style={-{Latex[length=2mm]}, thick, draw=gray!65}
]
\node[main] (sources) {1. Επιλογή επίσημων και τραπεζικών πηγών};
\node[main, below=0.28cm of sources] (snapshot) {2. Συλλογή, καθαρισμός και αποθήκευση στιγμιότυπου περιεχομένου};
\node[main, below=0.28cm of snapshot] (relevance) {3. Έλεγχος συνάφειας με LLM και απόρριψη μη σχετικών πηγών};
\node[main, below=0.28cm of relevance] (extract) {4. Εξαγωγή πεδίων σε κοινό JSON σχήμα};
\node[check, below=0.28cm of extract] (quality) {5. Δομικός και σημασιολογικός έλεγχος με URL, χρόνο και αρχικό κείμενο};
\node[check, below=0.28cm of quality] (admin) {6. Ανθρώπινη εποπτεία, διόρθωση, έγκριση και ανάδραση στη ροή};
\node[main, below=0.28cm of admin] (store) {7. Αποθήκευση εγκεκριμένης εγγραφής σε δομημένο JSON};
\node[main, below=0.28cm of store] (app) {8. Αξιοποίηση από διαδικτυακή εφαρμογή και μετρικές αξιολόγησης};

\draw[arrow] (sources) -- (snapshot);
\draw[arrow] (snapshot) -- (relevance);
\draw[arrow] (relevance) -- (extract);
\draw[arrow] (extract) -- (quality);
\draw[arrow] (quality) -- (admin);
\draw[arrow] (admin) -- (store);
\draw[arrow] (store) -- (app);
\end{tikzpicture}
